\section{Complete Practical Coding}
\label{sec:practical-16}

Keep the scientific structure stable and isolate temporary review formatting
behind one switch.

\begin{latexcode}{Complete practical file: journal-ready generic article}
\documentclass[11pt]{article}
\usepackage[a4paper,margin=25mm]{geometry}
\usepackage{setspace}

\newif\ifreviewcopy
\reviewcopytrue % change to \reviewcopyfalse for a compact draft
\ifreviewcopy
  \doublespacing
\fi

\title{Concise, Informative Article Title}
\author{First Author \and Second Author}
\date{}

\begin{document}
\maketitle

\begin{abstract}
State the question, method, main result, and conclusion.
\end{abstract}

\section{Introduction}
End with the objective or hypothesis.

\section{Materials and Methods}
Describe the design, measurements, and analysis.

\section{Results}
Report findings in the same order as the methods.

\section{Discussion}
Interpret the result, limitation, and implication.

\section{Conclusion}
Give the main answer without repeating the abstract.

\section*{Declarations}
\textbf{Competing interests:} The authors declare none.\par
\textbf{Data availability:} Replace with the accurate statement.\par
\textbf{Author contributions:} Replace with verified contributions.

\end{document}
\end{latexcode}

\begin{codelab}
Compile both review-switch states. When converting to a verified journal class,
move the scientific sections unchanged and replace only class-specific metadata
and declarations.
\end{codelab}
